\documentclass[11pt,a4paper,fontset=fandol]{ctexart}

\usepackage[a4paper,top=2.4cm,bottom=2.6cm,left=2.35cm,right=2.35cm,headheight=18pt,headsep=0.55cm,footskip=26pt]{geometry}
\usepackage{amsmath,amssymb,amsthm}
\usepackage{fancyhdr}
\usepackage{lastpage}
\usepackage{enumitem}
\usepackage{titlesec}
\usepackage{xcolor}
\usepackage{hyperref}
\usepackage{bookmark}
\usepackage[most]{tcolorbox}
\usepackage{setspace}
\usepackage{array}

\hypersetup{
  colorlinks=true,
  linkcolor=blue!50!black,
  urlcolor=blue!50!black,
  citecolor=blue!50!black
}

\setstretch{1.16}
\setlength{\parindent}{2em}
\setlength{\parskip}{0.35em}
\allowdisplaybreaks
\setlist[enumerate]{itemsep=0.32em, topsep=0.32em, leftmargin=2.3em}
\setlist[itemize]{itemsep=0.25em, topsep=0.25em, leftmargin=2em}
\setcounter{tocdepth}{2}

\definecolor{mainblue}{RGB}{36,83,140}
\definecolor{lightblue}{RGB}{241,246,252}
\definecolor{lightgray}{RGB}{246,246,246}
\definecolor{accent}{RGB}{153,76,0}
\definecolor{rulegray}{RGB}{110,118,128}
\definecolor{softgreen}{RGB}{236,247,240}

\titleformat{\section}
  {\Large\bfseries\color{mainblue}}
  {\thesection.}
  {0.45em}
  {}
  [\vspace{0.25em}\color{rulegray}\titlerule]
\titlespacing*{\section}{0pt}{1.3em}{0.9em}
\titleformat{\subsection}{\large\bfseries\color{mainblue}}{\thesubsection}{0.5em}{}
\titlespacing*{\subsection}{0pt}{1.05em}{0.55em}
\titleformat{\subsubsection}{\normalsize\bfseries\color{accent}}{\thesubsubsection}{0.5em}{}

\pagestyle{fancy}
\fancyhf{}
\fancyhead[L]{\small 组合专题课堂讲义}
\fancyhead[C]{\small 容斥原理入门与提高}
\fancyhead[R]{\small 刘文博}
\fancyfoot[C]{\small 第 \thepage 页 / 共 \pageref{LastPage} 页}
\renewcommand{\headrulewidth}{0.55pt}
\renewcommand{\footrulewidth}{0.35pt}

\newtheoremstyle{formaltheorem}
  {0.8em}{0.8em}{\normalfont}{}{\bfseries\color{mainblue}}{．}{0.6em}{}
\newtheoremstyle{formalexample}
  {0.8em}{0.8em}{\normalfont}{}{\bfseries\color{accent}}{．}{0.6em}{}

\theoremstyle{formaltheorem}
\newtheorem{theorem}{定理}[section]
\newtheorem{method}{方法}[section]
\theoremstyle{formalexample}
\newtheorem{example}{例题}[section]
\newtheorem{exercise}{练习}[section]

\newenvironment{solution}{\par\noindent\textbf{解：}}{\hfill$\square$\par}

\newtcolorbox{goalbox}[1]{
  breakable,
  colback=lightblue,
  colframe=mainblue,
  boxrule=0.7pt,
  arc=1mm,
  left=1.2mm,
  right=1.2mm,
  top=1mm,
  bottom=1mm,
  title=\textbf{#1},
  fonttitle=\bfseries
}

\newtcolorbox{notebox}[1]{
  breakable,
  colback=lightgray,
  colframe=black!50,
  boxrule=0.55pt,
  arc=1mm,
  left=1.2mm,
  right=1.2mm,
  top=1mm,
  bottom=1mm,
  title=\textbf{#1},
  fonttitle=\bfseries
}

\newtcolorbox{skillbox}[1]{
  breakable,
  colback=softgreen,
  colframe=green!40!black,
  boxrule=0.65pt,
  arc=1mm,
  left=1.2mm,
  right=1.2mm,
  top=1mm,
  bottom=1mm,
  title=\textbf{#1},
  fonttitle=\bfseries
}

\newcommand{\Dn}{D_n}

\begin{document}

\thispagestyle{empty}
\begin{titlepage}
\centering
\vspace*{0.9cm}

\begin{tcolorbox}[
  enhanced,
  width=0.92\textwidth,
  colback=white,
  colframe=mainblue,
  boxrule=1pt,
  arc=0mm,
  left=6mm,
  right=6mm,
  top=8mm,
  bottom=8mm
]
\centering

{\fontsize{24}{30}\selectfont\bfseries 组合专题课堂讲义\par}
\vspace{0.7cm}
{\fontsize{17}{22}\selectfont 容斥原理入门与提高\par}

\vspace{1.1cm}
\color{rulegray}\rule{0.72\textwidth}{0.7pt}\par
\vspace{1cm}

\begin{tcolorbox}[colback=lightblue,colframe=mainblue,width=0.82\textwidth,boxrule=1pt]
\centering
{\Large 适合对象：小学高年级至初中低年级数学竞赛学生}\\[0.4em]
{\large 已具备基础计数能力，准备学习“有重叠限制”的系统处理方法}
\end{tcolorbox}

\vspace{1.2cm}

\renewcommand{\arraystretch}{1.42}
\begin{tabular}{>{\bfseries}p{3.5cm}p{8cm}}
讲义作者： & 刘文博 \\
专题关键词： & 容斥原理、补集法、重叠计数、部分位置受限、错位排列 \\
建议课时： & 2--3 课时 \\
专题目标： & 学会处理“至少满足一个条件”“多个坏事件重叠”这两类典型计数问题 \\
编译方式： & XeLaTeX \\
\end{tabular}

\vspace{1.2cm}
\color{rulegray}\rule{0.72\textwidth}{0.7pt}\par
\vspace{0.8cm}

{\normalsize 本讲从“为什么会多算”出发，建立容斥原理的直观，再把它应用到倍数计数、部分位置受限、错位排列等竞赛常见模型中。\par}

\vspace{0.5cm}
{\large \today\par}
\end{tcolorbox}
\end{titlepage}

\pagenumbering{Roman}
\pagestyle{plain}
\tableofcontents
\newpage
\pagestyle{fancy}
\pagenumbering{arabic}

\section{专题目标与核心问题}

\begin{goalbox}{这一讲要解决什么问题}
前面学加法原理、乘法原理时，很多题目都能“分开算再相加”或“分步算再相乘”。
但当不同条件之间会\textbf{重叠}时，直接相加常常会\textbf{重复计算}。
容斥原理正是处理这类重复计数问题的标准工具。
\end{goalbox}

\begin{goalbox}{学完以后希望达到的能力}
\begin{itemize}
  \item 能说清楚为什么“直接相加”会多算；
  \item 会使用两个集合、三个集合的容斥公式；
  \item 能把“至少一个满足”“都不满足”转化为容斥问题；
  \item 能处理竞赛中常见的“部分位置受限”与“错位排列”题目。
\end{itemize}
\end{goalbox}

\begin{notebox}{本讲和补集法的关系}
补集法其实是容斥原理的一种最简单的表现：
\[
\text{满足要求}=\text{总数}-\text{不满足要求}.
\]
当“不满足要求”只有一类时，补集法就够了；当“不满足要求”分成很多互相重叠的坏事件时，就需要进一步使用容斥原理。
\end{notebox}

\section{为什么会出现“容斥”}

\subsection{从一个最简单的重叠开始}

\begin{example}
在 $1$ 到 $30$ 中，求“是 2 的倍数或 3 的倍数”的整数个数。
\end{example}

\begin{solution}
如果直接算：
\[
\lfloor 30/2\rfloor=15,
\qquad
\lfloor 30/3\rfloor=10.
\]
相加得到 $15+10=25$。

但这个答案偏大，因为既是 2 的倍数又是 3 的倍数，也就是 6 的倍数，被重复计算了两次。

而 1 到 30 中，6 的倍数有
\[
\lfloor 30/6\rfloor=5
\]
个。

所以正确答案应为
\[
15+10-5=20.
\]
这就是最朴素的“先加后减”的思想。
\end{solution}

\subsection{两个集合公式}

\begin{theorem}
若 $A,B$ 是两个有限集合，则
\[
|A\cup B|=|A|+|B|-|A\cap B|.
\]
\end{theorem}

\begin{solution}
把属于 $A\cap B$ 的元素想象成被同时算进了 $|A|$ 和 $|B|$，于是它们被重复计算了 2 次，而在并集里本来只应算 1 次，所以需要减去 1 次。
\end{solution}

\begin{skillbox}{看到这类词就要想到容斥}
\begin{itemize}
  \item “至少满足一个条件”
  \item “是 A 类或 B 类或 C 类”
  \item “不允许出现某些坏情况”
  \item “有若干个位置受限制”
\end{itemize}
\end{skillbox}

\section{三个集合与一般形式}

\subsection{三个集合公式}

\begin{theorem}
若 $A,B,C$ 是三个有限集合，则
\[
|A\cup B\cup C|
=|A|+|B|+|C|-|A\cap B|-|A\cap C|-|B\cap C|+|A\cap B\cap C|.
\]
\end{theorem}

\begin{solution}
先把三个集合的大小全部相加时，属于两个集合交集的元素被多算了，属于三个集合公共部分的元素被算了 3 次。

减去两两交集后，三者公共部分又被多减了，所以最后还要加回一次。于是符号呈现出
\[
+,-,+
\]
交替出现的规律。
\end{solution}

\begin{notebox}{三个集合为什么是“加减加”}
可以把它记成一句话：
\begin{itemize}
  \item 单个集合先全加；
  \item 两两交集要减掉；
  \item 三者公共部分再补回来。
\end{itemize}
以后推广到更多集合时，也是“奇数层加，偶数层减”的交替结构。
\end{notebox}

\subsection{一般形式先知道就够}
若坏事件有 $A_1,A_2,\ldots,A_n$，则
\[
\left|\bigcup_{i=1}^n A_i\right|
=\sum |A_i|-
\sum |A_i\cap A_j|+
\sum |A_i\cap A_j\cap A_k|-\cdots.
\]
对现阶段训练来说，不必死记一般公式，更重要的是理解：
\begin{itemize}
  \item 先加单个坏事件；
  \item 再减两两重叠；
  \item 再加三重重叠；
  \item 一直交替下去。
\end{itemize}

\section{计数题里的三种常见用法}

\subsection{用法一：求“至少满足一个条件”}

\begin{example}
在 $1$ 到 $100$ 中，是 2 的倍数或 3 的倍数或 5 的倍数的整数有多少个？
\end{example}

\begin{solution}
设
\begin{align*}
A&=\{\text{2 的倍数}\},\\
B&=\{\text{3 的倍数}\},\\
C&=\{\text{5 的倍数}\}.
\end{align*}

先算单个集合：
\[
|A|=50,\qquad |B|=33,\qquad |C|=20.
\]

再算两两交集：
\[
|A\cap B|=\lfloor 100/6\rfloor=16,
\]
\[
|A\cap C|=\lfloor 100/10\rfloor=10,
\qquad
|B\cap C|=\lfloor 100/15\rfloor=6.
\]

三者交集：
\[
|A\cap B\cap C|=\lfloor 100/30\rfloor=3.
\]

所以
\[
|A\cup B\cup C|=50+33+20-16-10-6+3=74.
\]

因此共有
\[
74
\]
个整数。
\end{solution}

\subsection{用法二：求“都不满足若干坏条件”}

\begin{example}
4 个人甲、乙、丙、丁排座位。要求甲不坐 1 号位，乙不坐 2 号位，丙不坐 3 号位，丁没有限制。共有多少种排法？
\end{example}

\begin{solution}
设
\[
A=\text{“甲坐 1 号位”},\quad B=\text{“乙坐 2 号位”},\quad C=\text{“丙坐 3 号位”}.
\]

我们真正要求的是：坏事件 $A,B,C$ 都不发生的排法数。

总排法有
\[
4!=24
\]
种。

单个坏事件：
\[
|A|=|B|=|C|=3!=6,
\]
所以单项总数为 $18$。

两两交集：
\[
|A\cap B|=|A\cap C|=|B\cap C|=2!=2,
\]
总数为 $6$。

三重交集：
\[
|A\cap B\cap C|=1!=1.
\]

因此满足要求的排法为
\[
24-18+6-1=11.
\]
\end{solution}

\subsection{用法三：处理部分位置受限与错位排列}

\begin{example}
6 个不同数字 $1,2,3,4,5,6$ 排成一行，要求 $1,2,3$ 都不在自己的位置上，而 $4,5,6$ 没有限制。共有多少种排法？
\end{example}

\begin{solution}
设
\[
A_1=\text{“1 在第 1 位”},\quad A_2=\text{“2 在第 2 位”},\quad A_3=\text{“3 在第 3 位”}.
\]

总排法有
\[
6!=720
\]
种。

单个坏事件共有
\[
3\times 5!=3\times 120=360
\]
种。

两个坏事件同时发生共有
\[
\binom32 4!=3\times 24=72
\]
种。

三个坏事件同时发生有
\[
3!=6
\]
种。

因此满足要求的排法数为
\[
720-360+72-6=426.
\]

这类题最容易误判成 $D_3$ 或 $D_6$，其实都不对，因为这里只限制了前 3 个对象。
\end{solution}

\section{错位排列是容斥原理的重要应用}

\begin{example}
用容斥原理计算 $D_5$。
\end{example}

\begin{solution}
设 $A_i$ 表示“第 $i$ 个对象在原位”，其中 $i=1,2,3,4,5$。

那么 $D_5$ 就等于总排法减去至少有一个对象在原位的排法：
\[
D_5=5!-\left|A_1\cup A_2\cup A_3\cup A_4\cup A_5\right|.
\]

按容斥展开：
\[
D_5=5!-\binom51 4!+\binom52 3!-\binom53 2!+\binom54 1!-\binom55 0!.
\]

逐项计算：
\[
5!=120,
\]
\[
\binom51 4!=5\times 24=120,
\]
\[
\binom52 3!=10\times 6=60,
\]
\[
\binom53 2!=10\times 2=20,
\]
\[
\binom54 1!=5,
\qquad
\binom55 0!=1.
\]

因此
\[
D_5=120-120+60-20+5-1=44.
\]

这和递推得到的结果完全一致。
\end{solution}

\begin{skillbox}{和错位排列结合时的固定套路}
若共有 $n$ 个对象，而其中前 $m$ 个对象不能回原位，则可以把
\[
A_i=\text{“第 $i$ 个受限制对象回到原位”}
\qquad (1\le i\le m)
\]
作为坏事件，最后写成
\[
\sum_{k=0}^{m}(-1)^k\binom{m}{k}(n-k)!
\]
的形式。
\end{skillbox}

\section{解题步骤与常见误区}

\begin{method}
遇到容斥题时，建议按下面 4 步走：
\begin{enumerate}
  \item 先确定“好事件”还是“坏事件”更容易描述；
  \item 给每一个坏事件取名字，如 $A_1,A_2,A_3$；
  \item 依次计算单项、两两交集、三重交集；
  \item 最后检查符号是否是“加减加减”交替。
\end{enumerate}
\end{method}

\begin{notebox}{最常见的 5 个错误}
\begin{itemize}
  \item 只把单个坏事件减掉，忘了把交集加回来；
  \item 交集个数会算，但忘了乘组合数 $\binom{m}{k}$；
  \item 把“只有部分对象受限”的题误判成标准错位排列；
  \item 恰好有 $n-1$ 个在原位这种不可能情况被误算进去；
  \item 计算完后不检查答案是否超过总数 $n!$。
\end{itemize}
\end{notebox}

\section{课堂练习}

\begin{exercise}
在 $1$ 到 $60$ 中，是 4 的倍数或 6 的倍数的整数有多少个？
\end{exercise}

\begin{exercise}
在 $1$ 到 $120$ 中，是 3 的倍数或 5 的倍数或 8 的倍数的整数有多少个？
\end{exercise}

\begin{exercise}
5 个人甲、乙、丙、丁、戊排座位。要求甲不坐 1 号位，乙不坐 2 号位，丙不坐 3 号位，其余两人没有限制。共有多少种排法？
\end{exercise}

\begin{exercise}
7 个对象排成一行，要求前 4 个对象都不在原位，后 3 个对象没有限制。共有多少种排法？
\end{exercise}

\begin{exercise}
用容斥原理求 $D_4$。
\end{exercise}

\begin{exercise}
8 个不同数字排成一行，要求至少有 1 个数字在原位。共有多少种排法？
\end{exercise}

\newpage
\appendix
\section{练习解答}

\textbf{1}\begin{solution}
设
\[
A=\{\text{4 的倍数}\},\qquad B=\{\text{6 的倍数}\}.
\]

则
\[
|A|=\lfloor 60/4\rfloor=15,
\qquad
|B|=\lfloor 60/6\rfloor=10.
\]

两者同时满足时，是 12 的倍数，所以
\[
|A\cap B|=\lfloor 60/12\rfloor=5.
\]

因此
\[
|A\cup B|=15+10-5=20.
\]

所以共有
\[
20
\]
个整数。
\end{solution}

\textbf{2}\begin{solution}
设 $A,B,C$ 分别表示“是 3 的倍数”“是 5 的倍数”“是 8 的倍数”。

单个集合：
\[
|A|=40,\qquad |B|=24,\qquad |C|=15.
\]

两两交集：
\[
|A\cap B|=\lfloor 120/15\rfloor=8,
\]
\[
|A\cap C|=\lfloor 120/24\rfloor=5,
\qquad
|B\cap C|=\lfloor 120/40\rfloor=3.
\]

三重交集：
\[
|A\cap B\cap C|=\lfloor 120/120\rfloor=1.
\]

由容斥原理：
\[
40+24+15-8-5-3+1=64.
\]

所以共有
\[
64
\]
个整数。
\end{solution}

\textbf{3}\begin{solution}
设
\[
A=\text{“甲坐 1 号位”},\quad B=\text{“乙坐 2 号位”},\quad C=\text{“丙坐 3 号位”}.
\]

总排法有
\[
5!=120
\]
种。

单个坏事件共有
\[
3\times 4!=3\times 24=72
\]
种。

两两交集共有
\[
\binom32 3!=3\times 6=18
\]
种。

三重交集有
\[
2!=2
\]
种。

所以满足要求的排法数为
\[
120-72+18-2=64.
\]
\end{solution}

\textbf{4}\begin{solution}
设前 4 个对象中，第 $i$ 个回原位的事件为 $A_i$。则满足要求的排列数为
\[
7!-\binom41 6!+\binom42 5!-\binom43 4!+\binom44 3!.
\]

逐项计算：
\[
7!=5040,
\qquad
\binom41 6!=4\times 720=2880,
\]
\[
\binom42 5!=6\times 120=720,
\qquad
\binom43 4!=4\times 24=96,
\]
\[
\binom44 3!=6.
\]

所以所求为
\[
5040-2880+720-96+6=2790.
\]
\end{solution}

\textbf{5}\begin{solution}
由容斥公式，
\[
D_4=4!-\binom41 3!+\binom42 2!-\binom43 1!+\binom44 0!.
\]

代入得
\[
D_4=24-4\times 6+6\times 2-4\times 1+1.
\]

即
\[
D_4=24-24+12-4+1=9.
\]
\end{solution}

\textbf{6}\begin{solution}
“至少有 1 个数字在原位”最适合用补集法。

总排法有
\[
8!=40320
\]
种。

没有任何数字在原位的排法数就是
\[
D_8.
\]
先由递推算得：
\[
D_7=1854,
\qquad
D_8=7(D_7+D_6)=7(1854+265)=14833.
\]

因此至少有 1 个数字在原位的排法数为
\[
8!-D_8=40320-14833=25487.
\]
\end{solution}

\end{document}
